
\section{\BookRecordIndex{T0340}稳定映射的\BookRecordIndex{T0139}格罗莫夫紧性}\label{sec:5-5}
\zhindex{wendingyingshedegeluomofujinxing}{稳定映射的格罗莫夫紧性}

本节先为（亏格为零的）\BookRecordIndex{T0340}稳定映射列定义\BookRecordIndex{T0141}格罗莫夫收敛，并证明每一列\BookRecordIndex{T0340}稳定映射都有一个
\BookRecordIndex{T0141}格罗莫夫收敛子列。然后证明极限在定义~\ref{def:5-1-4} 的等价关系下唯一。接下来，
我们用适当的非负数
$\rho_\varepsilon(\mathbf x,\mathbf x')\ge0$
\sourcefootnote{原书此处写作 $\rho_\varepsilon(\mathbf x,\mathbf x')>0$；但本节定义允许
其取零，且注~\ref{rem:5-5-7} 明确指出等价\BookRecordIndex{T0340}稳定映射之间的该量为零，故补为非负。}
衡量两个\BookRecordIndex{T0340}稳定映射
$\mathbf x=(\mathbf u,\mathbf z)$ 与
$\mathbf x'=(\mathbf u',\mathbf z')$ 之间的距离。这些数并不定义度量，但当
$\varepsilon$ 充分小时，它们具有如下性质：$\mathbf x^\nu$ 收敛到 $\mathbf x$，当且仅当
$\rho_\varepsilon(\mathbf x,\mathbf x^\nu)$ 在 $\nu\to\infty$ 时收敛到零。第~\ref{sec:5-6}节
将说明，这些结果会在\BookRecordIndex{T0340}稳定映射\BookRecordIndex{T0246}模空间上给出一个自然拓扑。

\paragraph{\BookRecordIndex{T0141}格罗莫夫收敛。}
设\BookRecordIndex{T0340}稳定映射 $(\mathbf u,\mathbf z)\in\mathcal{SC}_{0,n}(M;J)$ 以树 $T$ 为模型。
对顶点 $\alpha\in T$ 和开集 $U_\alpha\subset S^2$，记
\[
  E_\alpha(\mathbf u;U_\alpha)
  :=E(u_\alpha;U_\alpha)
    +\sum_{\substack{\beta\in T\\
                     \alpha E\beta,\ z_{\alpha\beta}\in U_\alpha}}
      m_{\alpha\beta}(\mathbf u).
\]
对 $z\in S^2$ 和 $\varepsilon>0$，用 $B_\varepsilon(z)\subset S^2$ 表示相对于
Fubini--Study 度量、中心为 $z$、半径为 $\varepsilon$ 的开球。

\begin{definition}[\BookRecordIndex{T0141}格罗莫夫收敛]\label{def:5-5-1}
设 $J^\nu\in\mathcal J_\tau(M,\omega)$ 是 $M$ 上一列被 $\omega$ \BookRecordIndex{T0372}驯服的\BookRecordIndex{T0088}殆复结构，
并收敛到 $J\in\mathcal J_\tau(M,\omega)$。设\BookRecordIndex{T0340}稳定映射列
\[
  (\mathbf u^\nu,\mathbf z^\nu)
  =\bigl(\{u_\alpha^\nu\}_{\alpha\in T^\nu},
    \{z_{\alpha\beta}^\nu\}_{\alpha E^\nu\beta},
    \{\alpha_i^\nu,z_i^\nu\}_{1\le i\le n}\bigr)
  \in\mathcal{SC}_{0,n}(M;J^\nu).
\]
若对每个充分大的 $\nu$，都存在满射\BookRecordIndex{T0305}树同态
\[
  f^\nu:T\longrightarrow T^\nu
\]
和一族 \BookRecordIndex{T0040}Möbius 变换 $\{\phi_\alpha^\nu\}_{\alpha\in T}$，使下列条件成立，就称
$(\mathbf u^\nu,\mathbf z^\nu)$ \emph{\BookRecordIndex{T0141}格罗莫夫收敛}到\BookRecordIndex{T0340}稳定映射
\[
  (\mathbf u,\mathbf z)
  =\bigl(\{u_\alpha\}_{\alpha\in T},
    \{z_{\alpha\beta}\}_{\alpha E\beta},
    \{\alpha_i,z_i\}_{1\le i\le n}\bigr)
  \in\mathcal{SC}_{0,n}(M;J).
\]

\medskip
\noindent\textup{（映射）}
对每个 $\alpha\in T$，映射列
\[
  u_{f^\nu(\alpha)}^\nu\circ\phi_\alpha^\nu:S^2\to M
\]
在 $S^2\setminus Z_\alpha$ 上\BookRecordIndex{T0416}逐紧一致收敛到 $u_\alpha$。

\medskip
\noindent\textup{（能量）}
若 $\alpha E\beta$，则
\begin{equation}\label{eq:5-5-1}
  m_{\alpha\beta}(\mathbf u)
  =\lim_{\varepsilon\to0}\lim_{\nu\to\infty}
  E_{f^\nu(\alpha)}
  \bigl(\mathbf u^\nu;
    \phi_\alpha^\nu(B_\varepsilon(z_{\alpha\beta}))\bigr).
\end{equation}

\medskip
\noindent\textup{（重标度）}
设 $\alpha,\beta\in T$ 满足 $\alpha E\beta$，并设 $\nu_j$ 是一个满足
$f^{\nu_j}(\alpha)=f^{\nu_j}(\beta)$ 的子列。则
\[
  \phi_{\alpha\beta}^{\nu_j}
  :=(\phi_\alpha^{\nu_j})^{-1}\circ\phi_\beta^{\nu_j}
\]
在 $S^2\setminus\{z_{\beta\alpha}\}$ 上\BookRecordIndex{T0416}逐紧一致收敛到 $z_{\alpha\beta}$。

\medskip
\noindent\textup{（\BookRecordIndex{T0191}节点）}
设 $\alpha,\beta\in T$ 满足 $\alpha E\beta$，并设 $\nu_j$ 是一个满足
$f^{\nu_j}(\alpha)\ne f^{\nu_j}(\beta)$ 的子列。则
\[
  z_{\alpha\beta}
  =\lim_{j\to\infty}
  (\phi_\alpha^{\nu_j})^{-1}
  \bigl(z_{f^{\nu_j}(\alpha)f^{\nu_j}(\beta)}^{\nu_j}\bigr).
\]

\medskip
\noindent\textup{（\BookRecordIndex{T0066}标记点）}
对所有 $i=1,\ldots,n$，有
\[
  \alpha_i^\nu=f^\nu(\alpha_i),
  \qquad
  z_i=\lim_{\nu\to\infty}
  (\phi_{\alpha_i}^\nu)^{-1}(z_i^\nu).
\]
\end{definition}

（重标度）、（\BookRecordIndex{T0191}节点）和（\BookRecordIndex{T0066}标记点）公理只涉及底\BookRecordIndex{T0072}层\BookRecordIndex{T0396}预稳定曲线 $\mathbf z^\nu$ 的极限行为，
与定义 D.5.1 中\BookRecordIndex{T0339}稳定曲线 \BookRecordIndex{T0022}GK 收敛的相应部分完全相同。（映射）和（能量）公理则是
定义~\ref{def:5-2-1} 中相应公理的自然推广。注意，固定 $T$ 后，能由 $T$ 满射到达的树
$T'$ 只有有限多个同构类型；因此，在模掉同构之后，$T^\nu$ 与 $f^\nu$ 只有有限种可能。
所以，一个\BookRecordIndex{T0141}格罗莫夫收敛列是有限多个子列的并，并且在每个子列上都可把 $T^\nu$ 和
$f^\nu$ 取为常值。关于\BookRecordIndex{T0340}稳定映射列的重要结果因而可以化归到已经证明的全纯球面列结果。
关键观察，是把定义~\ref{def:5-5-1} 中\BookRecordIndex{T0340}稳定映射的\BookRecordIndex{T0141}格罗莫夫收敛，与
定义~\ref{def:5-2-1} 中单个 $J$-全纯球面列的\BookRecordIndex{T0141}格罗莫夫收敛联系起来。

设 $(\mathbf u,\mathbf z)\in\mathcal{SC}_{0,n}(M;J)$ 是以\BookRecordIndex{T0084}带标号树
$(T,E,\Lambda)$ 为模型的\BookRecordIndex{T0340}稳定映射，并设 $T_0\subset T$ 是一个\BookRecordIndex{T0417}子树。
$(\mathbf u,\mathbf z)$ \emph{在 $T_0$ 上的限制}是\BookRecordIndex{T0340}稳定映射
$(\mathbf u_0,\mathbf z_0)$，其模型是\BookRecordIndex{T0084}带标号树 $(T_0,E_0,\Lambda_0)$，其中
\[
  E_0:=E\cap(T_0\times T_0),
  \qquad
  \Lambda_{0\alpha}:=\Lambda_\alpha
    \cup\{\beta\in T\setminus T_0\mid \alpha E\beta\},
  \quad \alpha\in T_0.
\]
对满足 $\alpha E\beta$ 的 $\alpha,\beta\in T_0$，定义
\[
  u_{0\alpha}:=u_\alpha,
  \qquad z_{0\alpha\beta}:=z_{\alpha\beta}.
\]
其\BookRecordIndex{T0066}标记点以集合
\[
  I_0:=\{i\in I\mid \alpha_i\in T_0\}
  \cup\{\beta\in T\setminus T_0\mid
    \text{存在 }\alpha\in T_0\text{ 使 }\alpha E\beta\}
\]
为指标，并定义为
\[
  (\alpha_{0i},z_{0i}):=(\alpha_i,z_i)
  \quad(i\in I,\ \alpha_i\in T_0),
\]
以及
\[
  (\alpha_{0\beta},z_{0\beta}):=(\alpha,z_{\alpha\beta})
  \quad(\alpha\in T_0,\ \beta\in T\setminus T_0,\ \alpha E\beta).
\]
下面的命题将在本节稍后证明。

\begin{proposition}\label{prop:5-5-2}
设 $(\mathbf u^\nu,\mathbf z^\nu)\in\mathcal{SC}_{0,n}(M;J^\nu)$ 是一列以\BookRecordIndex{T0084}带标号树
$(T',E',\Lambda')$ 为模型的\BookRecordIndex{T0340}稳定映射，并通过满射\BookRecordIndex{T0305}树同态 $f:T\to T'$ \BookRecordIndex{T0141}格罗莫夫收敛到
以 $(T,E,\Lambda)$ 为模型的\BookRecordIndex{T0340}稳定映射
$(\mathbf u,\mathbf z)\in\mathcal{SC}_{0,n}(M;J)$。则对每个 $\alpha'\in T'$，
带\BookRecordIndex{T0066}标记点的 $J^\nu$-全纯球面列 $(u_{\alpha'}^\nu,Y_{\alpha'}^\nu)$ 在
定义~\ref{def:5-2-1} 的意义下\BookRecordIndex{T0141}格罗莫夫收敛到
$(\mathbf u,\mathbf z)$ 在\BookRecordIndex{T0417}子树 $f^{-1}(\alpha')\subset T$ 上的限制。
\end{proposition}

\begin{theorem}[极限的唯一性]\label{thm:5-5-3}
设 $J^\nu\in\mathcal J_\tau(M,\omega)$ 是 $M$ 上一列被 $\omega$ \BookRecordIndex{T0372}驯服的\BookRecordIndex{T0088}殆复结构，
并在 $C^\infty$ 拓扑下收敛到 $J\in\mathcal J_\tau(M,\omega)$。设\BookRecordIndex{T0340}稳定映射列
$(\mathbf u^\nu,\mathbf z^\nu)\in\mathcal{SC}_{0,n}(M;J^\nu)$ \BookRecordIndex{T0141}格罗莫夫收敛到两个
\BookRecordIndex{T0340}稳定映射 $(\mathbf u,\mathbf z)$ 和 $(\widetilde{\mathbf u},\widetilde{\mathbf z})$。
则这两个\BookRecordIndex{T0340}稳定映射等价。
\end{theorem}

\begin{proof}
取子列后，由前面的讨论，可以假设所有\BookRecordIndex{T0340}稳定映射
$(\mathbf u^\nu,\mathbf z^\nu)$ 都以同一棵树 $T'$ 为模型，而且\BookRecordIndex{T0141}格罗莫夫收敛定义中的
满射\BookRecordIndex{T0305}树同态
\[
  f^\nu:T\to T',
  \qquad
  \widetilde f^\nu:\widetilde T\to T'
\]
分别等于固定映射 $f:T\to T'$ 与
$\widetilde f:\widetilde T\to T'$。在此假设下，每个顶点 $\alpha'\in T'$ 决定一列带有
\BookRecordIndex{T0066}标记点 $Y_{\alpha'}^\nu$ 的 $J^\nu$-全纯曲线 $u_{\alpha'}^\nu$；这些\BookRecordIndex{T0066}标记点由原有\BookRecordIndex{T0066}标记点
与\BookRecordIndex{T0191}节点 $z_{\alpha'\beta'}^\nu$ 组成。由命题~\ref{prop:5-5-2}，映射列
$(u_{\alpha'}^\nu,Y_{\alpha'}^\nu)$ 在定义~\ref{def:5-2-1} 的意义下既\BookRecordIndex{T0141}格罗莫夫收敛到
$(\mathbf u,\mathbf z)$ 在\BookRecordIndex{T0417}子树 $f^{-1}(\alpha')$ 上的限制，也\BookRecordIndex{T0141}格罗莫夫收敛到
$(\widetilde{\mathbf u},\widetilde{\mathbf z})$ 在\BookRecordIndex{T0417}子树
$\widetilde f^{-1}(\alpha')$ 上的限制。因此，由定理~\ref{thm:5-4-1}，分别以
$f^{-1}(\alpha')$ 与 $\widetilde f^{-1}(\alpha')$ 为模型的两条\BookRecordIndex{T0176}既约曲线等价。
再由定义~\ref{def:5-1-4}，$(\mathbf u,\mathbf z)$ 与
$(\widetilde{\mathbf u},\widetilde{\mathbf z})$ 等价。
\end{proof}

\begin{theorem}[同伦类的收敛]\label{thm:5-5-4}
设 $J^\nu\in\mathcal J_\tau(M,\omega)$ 是 $M$ 上一列被 $\omega$ \BookRecordIndex{T0372}驯服的\BookRecordIndex{T0088}殆复结构，
并在 $C^\infty$ 拓扑下收敛到 $J\in\mathcal J_\tau(M,\omega)$。设\BookRecordIndex{T0340}稳定映射列
$(\mathbf u^\nu,\mathbf z^\nu)\in\mathcal{SC}_{0,n}(M;J^\nu)$ \BookRecordIndex{T0141}格罗莫夫收敛到
\BookRecordIndex{T0340}稳定映射
$(\mathbf u,\mathbf z)\in\mathcal{SC}_{0,n}(M;J)$
\sourcefootnote{原书把极限\BookRecordIndex{T0340}稳定映射所属空间印作
$\mathcal{SC}_{0,n}(M;J^\nu)$；但 $J^\nu$ 随 $\nu$ 变化，而格罗莫夫极限按
定义~\ref{def:5-5-1} 是 $J$-全纯\BookRecordIndex{T0340}稳定映射，故补为
$\mathcal{SC}_{0,n}(M;J)$。}。则：
\begin{enumerate}[label=(\roman*),leftmargin=3em]
\item 若
$x^\nu\in\bigcup_{\alpha\in T^\nu}u_\alpha^\nu(S^2)$ 收敛到 $x\in M$，则
$x\in\bigcup_{\alpha\in T}u_\alpha(S^2)$；
\item 当 $\nu$ 充分大时，\BookRecordIndex{T0227}连通和
$\#_{\alpha\in T^\nu}u_\alpha^\nu$ 与
$\#_{\alpha\in T}u_\alpha$ 同伦。
\end{enumerate}
\end{theorem}

\begin{proof}
与定理~\ref{thm:5-5-3} 的证明一样，只须考虑所有\BookRecordIndex{T0340}稳定映射都以同一棵树 $T'$ 为模型，
且配有同一个从 $T$ 到 $T'$ 的满射\BookRecordIndex{T0305}树同态的情形。在这一情形，对各顶点
$\alpha'\in T'$ 分别应用命题~\ref{prop:5-5-2} 和定理~\ref{thm:5-2-2}，即得结论。
\end{proof}

\begin{theorem}[\BookRecordIndex{T0139}格罗莫夫紧性]\label{thm:5-5-5}
设 $J^\nu\in\mathcal J_\tau(M,\omega)$ 是 $M$ 上一列被 $\omega$ \BookRecordIndex{T0372}驯服的\BookRecordIndex{T0088}殆复结构，
并在 $C^\infty$ 拓扑下收敛到 $J\in\mathcal J_\tau(M,\omega)$。设
$(\mathbf u^\nu,\mathbf z^\nu)\in\mathcal{SC}_{0,n}(M;J^\nu)$ 是满足
\[
  \sup_\nu E(\mathbf u^\nu)<\infty
\]
的\BookRecordIndex{T0340}稳定映射列。则 $(\mathbf u^\nu,\mathbf z^\nu)$ 有一个\BookRecordIndex{T0141}格罗莫夫收敛子列。
\end{theorem}

\begin{proof}
令 $c:=\sup_\nu E(\mathbf u^\nu)$。由命题~\ref{prop:4-1-4} 证明中的平均值估计，
可以对收敛族 $J^\nu\to J$ 一致选取一个正的球面能量下界
\sourcefootnote{原书称命题~\ref{prop:4-1-4} 的最大阈值
$\hbar(M,\omega,J)$ 连续依赖于 $J$，并据此\BookRecordIndex{T0326}推出
$\inf_\nu\hbar(M,\omega,J^\nu)>0$；正文并未证明这个由可实现全纯球面能量的下确界
所定义函数的连续性，一般也不应期待这种连续性。此处只需局部一致的正下界，而
$J^\nu\to J$ 所成紧参数族上的平均值估计常数可一致选取，已足以得到该结论。}，因而
\[
  \inf_\nu\hbar(M,\omega,J^\nu)>0.
\]
所以，由习题~\ref{ex:5-1-2}，在同构意义下，与能量不超过 $c$ 的稳定
$J^\nu$-全纯曲线对应的树只有有限多棵。取子列后，可以假设各树 $T^\nu$ 都同构。
对每个 $\nu$ 选取一个同构
\[
  T':=T^1\longrightarrow T^\nu.
\]
现在由定理~\ref{thm:5-3-1}，存在一个仍记为
$(\mathbf u^\nu,\mathbf z^\nu)$ 的子列，使得对每个 $\alpha'\in T'$，带\BookRecordIndex{T0066}标记点的
$J^\nu$-全纯球面列 $(u_{\alpha'}^\nu,Y_{\alpha'}^\nu)$ \BookRecordIndex{T0141}格罗莫夫收敛到一个带\BookRecordIndex{T0066}标记点的
稳定 $J$-全纯映射。随后，对每对 $\alpha'E'\beta'$ 引入一条附加边，把这些极限曲线连接
成一棵树，即得结论。
\end{proof}

下面的引理既用于证明命题~\ref{prop:5-5-2}，也用于下一节中\BookRecordIndex{T0340}稳定映射之间距离的数值衡量。
它说明，（重标度）与（\BookRecordIndex{T0191}节点）公理中的假设 $\alpha E\beta$ 可以删去；同时还给出
（\BookRecordIndex{T0066}标记点）公理的自然推广。为简单起见，只陈述\BookRecordIndex{T0305}树同态
$f^\nu:T\to T^\nu$ 固定并等于 $f:T\to T'$ 的情形。需要时，对 $\alpha\in T$ 使用
记号 $\alpha':=f(\alpha)$。因为 $f$ 是\BookRecordIndex{T0084}带标号树的同态，它必须把 $\alpha_i$ 映到
$T'$ 中承载第 $i$ 个\BookRecordIndex{T0066}标记点的分支 $\alpha_i'$；所以这个记号\BookRecordIndex{T0348}相容，$\alpha_i'$ 的含义
没有歧义。此外，对 $\gamma\in T_{\alpha\beta}$，约定
\[
  z_{\alpha\gamma}:=z_{\alpha\beta}.
\]

\begin{lemma}\label{lem:5-5-6}
设 $T=(T,E,\Lambda)$ 与 $T'=(T',E',\Lambda')$ 是带 $n$ 个标号的树，且
$f:T\to T'$ 是满射\BookRecordIndex{T0305}树同态。设
$\{\phi_\alpha^\nu\}_{\alpha\in T,\,\nu\in\mathbb N}$ 是一族满足
定义~\ref{def:5-5-1} 中（重标度）、（\BookRecordIndex{T0191}节点）与（\BookRecordIndex{T0066}标记点）公理的 \BookRecordIndex{T0040}Möbius 变换。则：
\begin{enumerate}[label=(\roman*),leftmargin=3em]
\item 若 $\alpha,\beta\in T$ 满足 $\alpha\ne\beta$ 且
$f(\alpha)=f(\beta)$，则映射列
\[
  \phi_{\alpha\beta}^\nu
  :=(\phi_\alpha^\nu)^{-1}\circ\phi_\beta^\nu
\]
在 $S^2\setminus\{z_{\beta\alpha}\}$ 上\BookRecordIndex{T0416}逐紧一致收敛到 $z_{\alpha\beta}$；
\item 若 $\alpha,\beta\in T$ 满足 $f(\alpha)\ne f(\beta)$，则
\[
  z_{\alpha\beta}
  =\lim_{\nu\to\infty}
  (\phi_\alpha^\nu)^{-1}
  \bigl(z_{f(\alpha)f(\beta)}^\nu\bigr);
\]
\item 对每个 $\alpha\in T$ 和每个 $i\in\{1,\ldots,n\}$，
\[
  z_{\alpha i}
  =\lim_{\nu\to\infty}
  (\phi_\alpha^\nu)^{-1}
  \bigl(z_{f(\alpha)i}^\nu\bigr).
\]
\end{enumerate}
\end{lemma}

\begin{proof}
断言 (i) 由定理~\ref{thm:5-4-1} 证明的步骤 1 得到。

下面证明 (ii)。对 $\alpha\in T$ 使用记号 $\alpha':=f(\alpha)$。设
$\gamma_0,\ldots,\gamma_m$ 是从 $\alpha$ 到 $\beta$ 的\BookRecordIndex{T0064}边链。选取 $\ell<m$，使得
\[
  \alpha'=\gamma_i'\quad(i\le\ell),
  \qquad
  \alpha'\ne\gamma_i'\quad(i>\ell).
\]
若 $\ell=0$，则由定义~\ref{def:5-5-1} 的（\BookRecordIndex{T0191}节点）公理直接有
\sourcefootnote{原书未单独处理\BookRecordIndex{T0305}树同态不塌缩首条边时可能出现的 $\ell=0$，而仍写出
$z_{\gamma_0\gamma_\ell}$；这会成为未定义的 $z_{\gamma_0\gamma_0}$。此时（\BookRecordIndex{T0191}节点）公理
直接给出所需极限；以下原证明则适用于 $\ell\ge1$，引理结论不变。}
\[
  z_{\alpha\beta}=z_{\gamma_0\gamma_1}
  =\lim_{\nu\to\infty}
    (\phi_{\gamma_0}^\nu)^{-1}
    \bigl(z_{\gamma_0'\gamma_1'}^\nu\bigr)
  =\lim_{\nu\to\infty}
    (\phi_\alpha^\nu)^{-1}(z_{\alpha'\beta'}^\nu).
\]

以下设 $\ell\ge1$。于是
\[
  z_{\alpha\beta}=z_{\gamma_0\gamma_1}=z_{\gamma_0\gamma_\ell},
  \qquad
  z_{\alpha'\beta'}^\nu
  =z_{\gamma_\ell'\gamma_{\ell+1}'}^\nu.
\]
由定义~\ref{def:5-5-1} 的（\BookRecordIndex{T0191}节点）公理，映射列
\[
  (\phi_{\gamma_\ell}^\nu)^{-1}
  \bigl(z_{\gamma_\ell'\gamma_{\ell+1}'}^\nu\bigr)
\]
收敛到
$z_{\gamma_\ell\gamma_{\ell+1}}\ne z_{\gamma_\ell\gamma_0}$。再由 (i)，
\[
  \begin{aligned}
  z_{\alpha\beta}=z_{\gamma_0\gamma_\ell}
  &=\lim_{\nu\to\infty}
    (\phi_{\gamma_0}^\nu)^{-1}
    \bigl(z_{\gamma_\ell'\gamma_{\ell+1}'}^\nu\bigr)\\
  &=\lim_{\nu\to\infty}
    (\phi_\alpha^\nu)^{-1}(z_{\alpha'\beta'}^\nu).
  \end{aligned}
\]
这证明了 (ii)。

为证明 (iii)，先注意，若 $\alpha=\alpha_i$，则
$z_{\alpha i}=z_i$ 且 $z_{\alpha'i}^\nu=z_i^\nu$，所以结论正是
定义~\ref{def:5-5-1} 的（\BookRecordIndex{T0066}标记点）公理。

若 $\alpha\ne\alpha_i$ 但 $\alpha'=\alpha_i'$，令
\[
  \zeta^\nu:=(\phi_{\alpha_i}^\nu)^{-1}(z_i^\nu).
\]
由（\BookRecordIndex{T0066}标记点）公理，$\zeta^\nu\to z_i$；而\BookRecordIndex{T0320}特殊点两两不同给出
$z_i\ne z_{\alpha_i\alpha}$。因此，由 (i)，转移映射
$(\phi_\alpha^\nu)^{-1}\circ\phi_{\alpha_i}^\nu$ 可在移动点 $\zeta^\nu$ 处取极限，
从而
\[
  \begin{aligned}
  z_{\alpha i}=z_{\alpha\alpha_i}
  &=\lim_{\nu\to\infty}
    (\phi_\alpha^\nu)^{-1}\circ\phi_{\alpha_i}^\nu(\zeta^\nu)\\
  &=\lim_{\nu\to\infty}(\phi_\alpha^\nu)^{-1}(z_i^\nu).
  \end{aligned}
\]
这里用移动点论证替代了原书的点集混写
\sourcefootnote{原书把上述点值极限的中间项印作
$(\phi_\alpha^\nu)^{-1}\circ\phi_{\alpha_i}^\nu(B_\delta(z_i))$，这是一个集合，
与等式两侧的点类型不合。令
$\zeta^\nu=(\phi_{\alpha_i}^\nu)^{-1}(z_i^\nu)\to z_i$，再用 (i) 在
$z_i\ne z_{\alpha_i\alpha}$ 的邻域上一致收敛，即唯一得到正文中的点值等式。}。

最后，若 $\alpha\ne\alpha_i$ 且 $\alpha'\ne\alpha_i'$，则
$z_{\alpha i}=z_{\alpha\alpha_i}$ 且
$z_{\alpha'i}^\nu=z_{\alpha'\alpha_i'}^\nu$。由 (ii)，
$(\phi_\alpha^\nu)^{-1}(z_{\alpha'i}^\nu)$ 收敛到 $z_{\alpha i}$。
引理得证。
\end{proof}

\begin{proof}[命题~\ref{prop:5-5-2} 的证明]
令 $T_0:=f^{-1}(\alpha')$，并以 $(\mathbf u_0,\mathbf z_0)$ 表示
$(\mathbf u,\mathbf z)$ 在 $T_0$ 上的限制。我们要证明，映射列
$(u_{\alpha'}^\nu,Y_{\alpha'}^\nu)$ 与变换族
$\{\phi_\alpha^\nu\}_{\alpha\in T_0,\,\nu\in\mathbb N}$ 满足
定义~\ref{def:5-2-1} 的要求。其（重标度）和（\BookRecordIndex{T0066}标记点）公理直接来自
定义~\ref{def:5-5-1} 中映射列 $(\mathbf u^\nu,\mathbf z^\nu)$ 的
（重标度）、（\BookRecordIndex{T0191}节点）与（\BookRecordIndex{T0066}标记点）公理。

为建立（映射）和（能量）公理，先由引理~\ref{lem:5-5-6} 得到
\begin{equation}\label{eq:5-5-2}
  \gamma\in T_{\alpha\beta}
  \quad\Longleftrightarrow\quad
  z_{\alpha\beta}
  =\lim_{\nu\to\infty}
  (\phi_\alpha^\nu)^{-1}(z_{\alpha'\gamma'}^\nu),
\end{equation}
其中 $\alpha,\beta,\gamma\in T$、$\alpha E\beta$，且
$\alpha':=f(\alpha)\ne\gamma':=f(\gamma)$。现在考虑集合
\[
  Z_{0\alpha}:=\{z_{\alpha\beta}\mid
    \beta\in T_0,\ \alpha E\beta\}.
\]
除非 $T=T_0$，否则对某些 $\alpha\in T_0$，它是 $Z_\alpha$ 的真子集。我们必须证明，
对每个 $\alpha\in T_0$，映射列
$u_{\alpha'}^\nu\circ\phi_\alpha^\nu$ 在
$S^2\setminus Z_{0\alpha}$ 上\BookRecordIndex{T0416}逐紧一致收敛到 $u_\alpha$。

设 $\alpha\in T_0$、$\beta\in T\setminus T_0$，且 $\alpha E\beta$。记
$\beta':=f(\beta)$。由式~\eqref{eq:5-5-2} 和定义~\ref{def:5-5-1} 的（能量）公理，
\[
  m_{\alpha\beta}(\mathbf u)
  =\lim_{\varepsilon\to0}\lim_{\nu\to\infty}
  \left(
    E\bigl(u_{\alpha'}^\nu\circ\phi_\alpha^\nu;
      B_\varepsilon(z_{\alpha\beta})\bigr)
    +m_{\alpha'\beta'}(\mathbf u^\nu)
  \right)
\]
以及
\[
  m_{\beta\alpha}(\mathbf u)
  =\lim_{\varepsilon\to0}\lim_{\nu\to\infty}
  \left(
    E\bigl(u_{\beta'}^\nu\circ\phi_\beta^\nu;
      B_\varepsilon(z_{\beta\alpha})\bigr)
    +m_{\beta'\alpha'}(\mathbf u^\nu)
  \right).
\]
另一方面，定义~\ref{def:5-5-1} 的（能量）与（映射）公理给出
\[
  \begin{aligned}
  E(\mathbf u)
  &=E(u_\alpha)+\sum_{\substack{\gamma\in T\\\alpha E\gamma}}
    m_{\alpha\gamma}(\mathbf u)\\
  &=\lim_{\varepsilon\to0}\lim_{\nu\to\infty}
    \Biggl(
      E\Bigl(u_{\alpha'}^\nu\circ\phi_\alpha^\nu;
        S^2\setminus\!\bigcup_{\substack{\gamma\in T\\\alpha E\gamma}}
        B_\varepsilon(z_{\alpha\gamma})\Bigr)\\[-0.4ex]
  &\hspace{8em}
      +\sum_{\substack{\gamma\in T\\\alpha E\gamma}}
       E_{\alpha'}\bigl(\mathbf u^\nu;
        \phi_\alpha^\nu(B_\varepsilon(z_{\alpha\gamma}))\bigr)
    \Biggr)\\
  &=\lim_{\nu\to\infty}
    \left(
      E(u_{\alpha'}^\nu\circ\phi_\alpha^\nu)
      +\sum_{\substack{\gamma'\in T'\\\alpha'E'\gamma'}}
       m_{\alpha'\gamma'}(\mathbf u^\nu)
    \right)
   =\lim_{\nu\to\infty}E(\mathbf u^\nu).
  \end{aligned}
\]
由于
\[
  E(\mathbf u)
  =m_{\alpha\beta}(\mathbf u)+m_{\beta\alpha}(\mathbf u),
  \qquad
  E(\mathbf u^\nu)
  =m_{\alpha'\beta'}(\mathbf u^\nu)
   +m_{\beta'\alpha'}(\mathbf u^\nu),
\]
可得
\begin{equation}\label{eq:5-5-3}
  \begin{gathered}
  m_{\alpha\beta}(\mathbf u)
  =\lim_{\nu\to\infty}m_{\alpha'\beta'}(\mathbf u^\nu),\\
  \lim_{\varepsilon\to0}\lim_{\nu\to\infty}
  E\bigl(u_{\alpha'}^\nu\circ\phi_\alpha^\nu;
    B_\varepsilon(z_{\alpha\beta})\bigr)=0.
  \end{gathered}
\end{equation}
这说明，对每个满足 $\beta\in T\setminus T_0$、$\alpha E\beta$ 的 $\beta$，映射列
$u_{\alpha'}^\nu\circ\phi_\alpha^\nu$ 在点 $z_{\alpha\beta}$ 附近不发生\BookRecordIndex{T0243}冒泡。因此，
它在 $S^2\setminus Z_{0\alpha}$ 上\BookRecordIndex{T0416}逐紧一致收敛到 $u_\alpha$。

最后验证定义~\ref{def:5-2-1} 的（能量）公理。设
$\alpha,\beta\in T_0$ 且 $\alpha E\beta$。由式~\eqref{eq:5-5-2} 与
式~\eqref{eq:5-5-3}，
\[
  \begin{aligned}
  m_{\alpha\beta}(\mathbf u_0)
  &=m_{\alpha\beta}(\mathbf u)
    -\sum_{\substack{\gamma_0,\gamma_1\in T_{\alpha\beta}\\
                      \gamma_0\in T_0,\ \gamma_1\notin T_0\\
                      \gamma_0E\gamma_1}}
      m_{\gamma_0\gamma_1}(\mathbf u)\\
  &=\lim_{\varepsilon\to0}\lim_{\nu\to\infty}
    \Biggl(
      E_{\alpha'}\bigl(\mathbf u^\nu;
        \phi_\alpha^\nu(B_\varepsilon(z_{\alpha\beta}))\bigr)
      -\sum_{\substack{\gamma\in T_{\alpha\beta}\\
                        \alpha'E'\gamma'}}
        m_{\alpha'\gamma'}(\mathbf u^\nu)
    \Biggr)\\
  &=\lim_{\varepsilon\to0}\lim_{\nu\to\infty}
    E\bigl(u_{\alpha'}^\nu\circ\phi_\alpha^\nu;
      B_\varepsilon(z_{\alpha\beta})\bigr).
  \end{aligned}
\]
命题得证。
\end{proof}

\paragraph{\BookRecordIndex{T0340}稳定映射之间距离的衡量。}
现在定义从一个\BookRecordIndex{T0340}稳定映射 $\mathbf x$ 到另一个\BookRecordIndex{T0340}稳定映射 $\mathbf x'$ 的“距离”
$\rho_\varepsilon(\mathbf x,\mathbf x')$。尽管有这个名称，它并非
$\mathbf x,\mathbf x'$ 的对称函数，也不满足三角不等式。设
\[
  \mathbf x=(\mathbf u,\mathbf z)
  =\bigl(\{u_\alpha\}_{\alpha\in T},
    \{z_{\alpha\beta}\}_{\alpha E\beta},
    \{\alpha_i,z_i\}_{1\le i\le n}\bigr)
  \in\mathcal{SC}_{0,n}(M,A;J)
\]
是带 $n$ 个\BookRecordIndex{T0066}标记点的亏格零稳定 \BookRecordIndex{T0003}$J$-全纯曲线，并固定充分小的常数
$\varepsilon>0$。对任意另一个\BookRecordIndex{T0340}稳定映射
\[
  \mathbf x'=(\mathbf u',\mathbf z')
  =\bigl(\{u_{\alpha'}'\}_{\alpha'\in T'},
    \{z_{\alpha'\beta'}'\}_{\alpha'E'\beta'},
    \{\alpha_i',z_i'\}_{1\le i\le n}\bigr)
  \in\mathcal{SC}_{0,n}(M,A;J),
\]
定义取值于 $[0,\infty]$ 的扩展非负实数
\[
  \rho_\varepsilon(\mathbf x,\mathbf x')
  :=\inf_{f:T\to T'}\ \inf_{\{\phi_\alpha\}}
  \rho_\varepsilon
  \bigl(\mathbf x,\mathbf x';f,\{\phi_\alpha\}\bigr),
\]
其中
\[
  \begin{aligned}
  &\rho_\varepsilon
    \bigl(\mathbf x,\mathbf x';f,\{\phi_\alpha\}\bigr)\\
  :={}&\sup_{\alpha E\beta}
    \Bigl|E_\alpha\bigl(\mathbf u;B_\varepsilon(z_{\alpha\beta})\bigr)
      -E_{f(\alpha)}\bigl(
        \mathbf u';\phi_\alpha(B_\varepsilon(z_{\alpha\beta}))\bigr)\Bigr|\\
  &+\sup_{\alpha\in T}\ 
    \sup_{S^2\setminus B_\varepsilon(Z_\alpha)}
    d\bigl(u_{f(\alpha)}'\circ\phi_\alpha,u_\alpha\bigr)\\
  &+\sup_{\substack{\alpha\ne\beta\\f(\alpha)=f(\beta)}}\ 
    \sup_{S^2\setminus B_\varepsilon(z_{\alpha\beta})}
    d\bigl(\phi_\beta^{-1}\circ\phi_\alpha,z_{\beta\alpha}\bigr)\\
  &+\sup_{f(\alpha)\ne f(\beta)}
    d\bigl(\phi_\beta^{-1}(z_{f(\beta)f(\alpha)}'),z_{\beta\alpha}\bigr)\\
  &+\sup_{\substack{\alpha\in T\\1\le i\le n}}
    d\bigl(\phi_\alpha^{-1}(z_{f(\alpha)i}'),z_{\alpha i}\bigr).
  \end{aligned}
\]
在 $\rho_\varepsilon$ 的定义中，下确界取遍所有
$\{\phi_\alpha\}_{\alpha\in T}\in G^T$，以及所有满足
$f(\alpha_i)=\alpha_i'$ 的满射\BookRecordIndex{T0305}树同态 $f:T\to T'$。若不存在这样的满射\BookRecordIndex{T0305}树同态，则规定
\[
  \rho_\varepsilon(\mathbf x,\mathbf x'):=\infty.
\]
上面的公式看起来颇为可畏，不过和式中的每一项都自然衡量了两个\BookRecordIndex{T0340}稳定映射可能不同的
一种方式。

\begin{remark}\label{rem:5-5-7}
依定义，函数
$\mathbf x'\mapsto\rho_\varepsilon(\mathbf x,\mathbf x')$ 可下降到\BookRecordIndex{T0340}稳定映射等价类的
\BookRecordIndex{T0246}模空间 $\overline{\mathcal M}_{0,n}(A;J)$ 上。因此，
\[
  \mathbf x'\equiv\mathbf y'
  \quad\Longrightarrow\quad
  \rho_\varepsilon(\mathbf x,\mathbf x')
  =\rho_\varepsilon(\mathbf x,\mathbf y'),
\]
而且
\[
  \mathbf x\equiv\mathbf x'
  \quad\Longrightarrow\quad
  \rho_\varepsilon(\mathbf x,\mathbf x')=0
\]
对所有
$\mathbf x,\mathbf x',\mathbf y'\in\mathcal{SC}_{0,n}(M,A;J)$ 成立。这里
$\equiv$ 表示范畴 $\mathcal{SC}_{0,n}(M,A;J)$ 上的等价关系。
\end{remark}

函数 $\rho_\varepsilon(\mathbf x,\mathbf x')$ 在如下意义下刻画\BookRecordIndex{T0141}格罗莫夫收敛。

\begin{lemma}\label{lem:5-5-8}
固定一条稳定 \BookRecordIndex{T0003}$J$-全纯曲线 $\mathbf x=(\mathbf u,\mathbf z)$。则存在常数
$\varepsilon_0>0$，使得对 $0<\varepsilon<\varepsilon_0$，下述结论成立：映射列
$\mathbf x^\nu=(\mathbf u^\nu,\mathbf z^\nu)$ \BookRecordIndex{T0141}格罗莫夫收敛到 $\mathbf x$，当且仅当
扩展非负实数列 $\rho_\varepsilon(\mathbf x,\mathbf x^\nu)$ 收敛到零。
\end{lemma}

\begin{proof}
选取充分小的 $\varepsilon_0>0$，使得对所有 $\alpha\in T$，
\[
  E\bigl(u_\alpha;B_{\varepsilon_0}(Z_\alpha)\bigr)<\hbar/2,
\]
并且
\[
  \alpha E\beta,\quad \alpha E\gamma,\quad\beta\ne\gamma
  \quad\Longrightarrow\quad
  B_{\varepsilon_0}(z_{\alpha\beta})
  \cap B_{\varepsilon_0}(z_{\alpha\gamma})=\varnothing.
\]
于是，对每个满足 $\alpha E\beta$ 的 $\alpha,\beta\in T$，都有
\[
  E_\alpha\bigl(\mathbf u;B_{\varepsilon_0}(z_{\alpha\beta})\bigr)
  =E\bigl(u_\alpha;B_{\varepsilon_0}(z_{\alpha\beta})\bigr)
   +m_{\alpha\beta}(\mathbf u).
\]
我们断言，这一 $\varepsilon_0$ 满足引理的要求。

若 $\mathbf x^\nu$ \BookRecordIndex{T0141}格罗莫夫收敛到 $\mathbf x$，则由定义和
引理~\ref{lem:5-5-6}，$\rho_\varepsilon(\mathbf x,\mathbf x^\nu)$ 收敛到零。
反过来，假设
\[
  \lim_{\nu\to\infty}
  \rho_\varepsilon(\mathbf x,\mathbf x^\nu)=0.
\]
当 $\nu\in\mathbb N$ 充分大时，存在满射\BookRecordIndex{T0305}树同态 $f^\nu:T\to T^\nu$ 和
$\{\phi_\alpha^\nu\}_{\alpha\in T}\in G^T$，使得
$f^\nu(\alpha_i)=\alpha_i^\nu$，且
\[
  \rho^\nu
  :=\rho_\varepsilon
    \bigl(\mathbf x,\mathbf x^\nu;f^\nu,
      \{\phi_\alpha^\nu\}\bigr)
  \le\rho_\varepsilon(\mathbf x,\mathbf x^\nu)+2^{-\nu}.
\]
我们证明这组数据满足定义~\ref{def:5-5-1} 的全部五个条件。以 $T$ 为定义域的满射\BookRecordIndex{T0305}树同态
只有有限种可能，所以只须考虑对所有 $\nu$ 都有
\[
  T^\nu=T^1=:T',
  \qquad
  f^\nu=f^1=:f:T\to T'
\]
的情形。需要时记 $\alpha':=f(\alpha)$。

（\BookRecordIndex{T0066}标记点）和（\BookRecordIndex{T0191}节点）公理显然成立。若 $\alpha E\beta$ 且 $\alpha'=\beta'$，则
$\rho^\nu\to0$ 说明
$(\phi_\alpha^\nu)^{-1}\circ\phi_\beta^\nu$ 在
$S^2\setminus B_\varepsilon(z_{\beta\alpha})$ 上收敛到 $z_{\alpha\beta}$，而
$(\phi_\beta^\nu)^{-1}\circ\phi_\alpha^\nu$ 在
$S^2\setminus B_\varepsilon(z_{\alpha\beta})$ 上收敛到 $z_{\beta\alpha}$。
由引理 D.1.4，前一个映射列在
$S^2\setminus\{z_{\beta\alpha}\}$ 上\BookRecordIndex{T0416}逐紧一致收敛到 $z_{\alpha\beta}$。
这证明了（重标度）公理。

为证明（映射）和（能量）公理，注意，由于 $\rho^\nu\to0$，映射列
\[
  u_\alpha^\nu:=u_{\alpha'}^\nu\circ\phi_\alpha^\nu
\]
在
$S^2\setminus\bigcup_{\alpha E\beta}B_\varepsilon(z_{\alpha\beta})$ 上一致收敛到
$u_\alpha$，而且当 $\alpha E\beta$ 时，
\begin{equation}\label{eq:5-5-4}
  m_{\alpha\beta}(\mathbf u)
  +E\bigl(u_\alpha;B_\varepsilon(z_{\alpha\beta})\bigr)
  =\lim_{\nu\to\infty}
  E_{\alpha'}\bigl(\mathbf u^\nu;
    \phi_\alpha^\nu(B_\varepsilon(z_{\alpha\beta}))\bigr).
\end{equation}
因此，由引理~\ref{lem:4-6-6}，$u_\alpha^\nu$ 在
$S^2\setminus\bigcup_{\alpha E\beta}B_\varepsilon(z_{\alpha\beta})$ 的紧子集上
一致且连同所有阶导数一起收敛到 $u_\alpha$。

还须证明，对任意 $\delta>0$，映射列 $u_\alpha^\nu$ 在
$B_\varepsilon(z_{\alpha\beta})\setminus B_\delta(z_{\alpha\beta})$ 中没有\BookRecordIndex{T0275}气泡。
由引理~\ref{lem:5-5-6}(ii)，对每个满足 $\alpha'\ne\gamma'$ 的 $\gamma\in T$，
\[
  z_{\alpha\gamma}
  =\lim_{\nu\to\infty}
  (\phi_\alpha^\nu)^{-1}(z_{\alpha'\gamma'}^\nu).
\]
这说明，对每个 $\delta>0$，存在 $\nu_0=\nu_0(\delta)$，使得当
$\nu\ge\nu_0$ 时，
\[
  \begin{aligned}
  &E_{\alpha'}\bigl(\mathbf u^\nu;
    \phi_\alpha^\nu(
      B_\varepsilon(z_{\alpha\beta})
      \setminus B_\delta(z_{\alpha\beta}))\bigr)\\
  &\qquad
  =E\bigl(u_\alpha^\nu;
      B_\varepsilon(z_{\alpha\beta})
      \setminus B_\delta(z_{\alpha\beta})\bigr).
  \end{aligned}
\]
所以，由式~\eqref{eq:5-5-4}，当 $0<\delta<\varepsilon$ 时，
\[
  \begin{aligned}
  &E\bigl(u_\alpha;
    B_\varepsilon(z_{\alpha\beta})
    \setminus B_\delta(z_{\alpha\beta})\bigr)\\
  &\qquad
  =\lim_{\nu\to\infty}
   E\bigl(u_\alpha^\nu;
    B_\varepsilon(z_{\alpha\beta})
    \setminus B_\delta(z_{\alpha\beta})\bigr).
  \end{aligned}
\]
由 $\varepsilon$ 的选取，这个极限不超过 $\hbar/2$，所以\BookRecordIndex{T0170}环域
$B_\varepsilon(z_{\alpha\beta})\setminus B_\delta(z_{\alpha\beta})$ 中不会发生\BookRecordIndex{T0243}冒泡。
因此，$u_\alpha^\nu$ 的每个子列都有一个进一步的子列，在
$B_\varepsilon(z_{\alpha\beta})\setminus\{z_{\alpha\beta}\}$ 上\BookRecordIndex{T0416}逐紧一致收敛到某条
\BookRecordIndex{T0003}$J$-全纯曲线。由\BookRecordIndex{T0333}唯一延拓定理，极限曲线总是 $u_\alpha$；所以整个映射列
$u_\alpha^\nu$ 都在该穿孔球上\BookRecordIndex{T0416}逐紧一致收敛到 $u_\alpha$。这证明了（映射）公理，
（能量）公理随即由式~\eqref{eq:5-5-4} 得到。因此
$\mathbf x^\nu=(\mathbf u^\nu,\mathbf z^\nu)$ \BookRecordIndex{T0141}格罗莫夫收敛到
$\mathbf x=(\mathbf u,\mathbf z)$，引理得证。
\end{proof}

函数 $\rho_\varepsilon$ 满足下面这个三角不等式的替代结论。

\begin{lemma}\label{lem:5-5-9}
设 $\mathbf x\in\mathcal{SC}_{0,n}(M,A;J)$，且
$0<\varepsilon<\varepsilon_0$，其中 $\varepsilon_0$ 如引理~\ref{lem:5-5-8} 中所选。
若 $\mathbf x'\in\mathcal{SC}_{0,n}(M,A;J)$ 满足
$\rho_\varepsilon(\mathbf x,\mathbf x')<\varepsilon$，且映射列
$\mathbf x^\nu\in\mathcal{SC}_{0,n}(M,A;J)$ \BookRecordIndex{T0141}格罗莫夫收敛到 $\mathbf x'$，则
\[
  \limsup_{\nu\to\infty}
  \rho_\varepsilon(\mathbf x,\mathbf x^\nu)
  \le\rho_\varepsilon(\mathbf x,\mathbf x').
\]
\end{lemma}

\begin{proof}
由假设，存在满射\BookRecordIndex{T0305}树同态 $g:T\to T'$ 和一族 \BookRecordIndex{T0040}Möbius 变换
$\{\psi_\alpha\}_{\alpha\in T}$，使得 $g(\alpha_i)=\alpha_i'$ 且
\begin{equation}\label{eq:5-5-5}
  \rho_\varepsilon
  \bigl(\mathbf x,\mathbf x';g,\{\psi_\alpha\}\bigr)<\varepsilon.
\end{equation}
需要时记 $\alpha':=g(\alpha)$。由式~\eqref{eq:5-5-5} 和
$\rho_\varepsilon$ 的定义，对每个满足 $\alpha'\ne\beta'$ 的
$\alpha,\beta\in T$ 以及每个 $i\in\{1,\ldots,n\}$，都有
\[
  d\bigl(\psi_\alpha^{-1}(z_{\alpha'\beta'}'),z_{\alpha\beta}\bigr)
  <\varepsilon,
  \qquad
  d\bigl(\psi_\alpha^{-1}(z_{\alpha'i}'),z_{\alpha i}\bigr)
  <\varepsilon.
\]
因此，对每个 $\alpha\in T$，
\begin{equation}\label{eq:5-5-6}
  Z_{\alpha'}(\mathbf x')
  \subset\psi_\alpha\bigl(B_\varepsilon(Z_\alpha(\mathbf x))\bigr).
\end{equation}

现设 $\mathbf x^\nu$ 通过满射\BookRecordIndex{T0305}树同态 $f^\nu:T'\to T^\nu$ 和 \BookRecordIndex{T0040}Möbius 变换
$\phi_{\alpha'}^\nu\in G$ \BookRecordIndex{T0141}格罗莫夫收敛到 $\mathbf x'$。我们断言
\begin{equation}\label{eq:5-5-7}
  \begin{aligned}
  &\rho_\varepsilon
  \bigl(\mathbf x,\mathbf x';g,\{\psi_\alpha\}\bigr)\\
  &\quad=\lim_{\nu\to\infty}
  \rho_\varepsilon\bigl(
    \mathbf x,\mathbf x^\nu;
    f^\nu\circ g,
    \{\phi_{g(\alpha)}^\nu\circ\psi_\alpha\}\bigr).
  \end{aligned}
\end{equation}
仍只须考虑对所有 $\nu$ 都有
\[
  T^\nu=T^1=:T'',
  \qquad
  f^\nu=f^1=:f:T'\to T''
\]
的情形。对 $\alpha\in T$ 简记
\[
  \alpha'':=f\circ g(\alpha),
  \qquad
  \chi_\alpha^\nu:=\phi_{\alpha'}^\nu\circ\psi_\alpha\in G,
  \qquad
  u_\alpha^\nu:=u_{\alpha''}^\nu\circ\chi_\alpha^\nu,
  \qquad
  Z_\alpha:=Z_\alpha(\mathbf x).
\]
由式~\eqref{eq:5-5-6}，映射列
$u_{\alpha''}^\nu\circ\phi_{\alpha'}^\nu$ 在
$S^2\setminus\psi_\alpha(B_\varepsilon(Z_\alpha))$ 上一致收敛到
$u_{\alpha'}'$。因而 $u_\alpha^\nu$ 在
$S^2\setminus B_\varepsilon(Z_\alpha)$ 上一致收敛到
$u_{\alpha'}'\circ\psi_\alpha$，所以
\[
  \sup_{S^2\setminus B_\varepsilon(Z_\alpha)}
  d\bigl(u_{\alpha'}'\circ\psi_\alpha,u_\alpha\bigr)
  =\lim_{\nu\to\infty}
  \sup_{S^2\setminus B_\varepsilon(Z_\alpha)}
  d(u_\alpha^\nu,u_\alpha).
\]

若 $\alpha,\beta\in T$ 满足 $\alpha\ne\beta$ 且
$\alpha'=\beta'$，则
\[
  \psi_\beta^{-1}\circ\psi_\alpha
  =(\chi_\beta^\nu)^{-1}\circ\chi_\alpha^\nu,
\]
从而
\[
  \sup_{S^2\setminus B_\varepsilon(z_{\alpha\beta})}
  d\bigl(\psi_\beta^{-1}\circ\psi_\alpha,z_{\beta\alpha}\bigr)
  =\sup_{S^2\setminus B_\varepsilon(z_{\alpha\beta})}
  d\bigl((\chi_\beta^\nu)^{-1}\circ\chi_\alpha^\nu,
    z_{\beta\alpha}\bigr).
\]

若 $\alpha'\ne\beta'$ 而 $\alpha''=\beta''$，则由
引理~\ref{lem:5-5-6}(i) 与式~\eqref{eq:5-5-6}，映射列
$(\phi_{\beta'}^\nu)^{-1}\circ\phi_{\alpha'}^\nu$ 在
$S^2\setminus\psi_\alpha(B_\varepsilon(z_{\alpha\beta}))$ 上一致收敛到
$z_{\beta'\alpha'}'$。所以
$(\chi_\beta^\nu)^{-1}\circ\chi_\alpha^\nu$ 在
$S^2\setminus B_\varepsilon(z_{\alpha\beta})$ 上一致收敛到
$\psi_\beta^{-1}(z_{\beta'\alpha'}')$，于是
\[
  \begin{aligned}
  d\bigl(\psi_\beta^{-1}(z_{\beta'\alpha'}'),z_{\beta\alpha}\bigr)
  =\lim_{\nu\to\infty}
  \sup_{S^2\setminus B_\varepsilon(z_{\alpha\beta})}
  d\bigl((\chi_\beta^\nu)^{-1}\circ\chi_\alpha^\nu,
    z_{\beta\alpha}\bigr).
  \end{aligned}
\]

若 $\alpha''\ne\beta''$，则
$(\phi_{\beta'}^\nu)^{-1}(z_{\beta''\alpha''}^\nu)$ 收敛到
$z_{\beta'\alpha'}'$，因而
\[
  d\bigl(\psi_\beta^{-1}(z_{\beta'\alpha'}'),z_{\beta\alpha}\bigr)
  =\lim_{\nu\to\infty}
  d\bigl((\chi_\beta^\nu)^{-1}(z_{\beta''\alpha''}^\nu),
    z_{\beta\alpha}\bigr).
\]
由引理~\ref{lem:5-5-6}(iii)，对每个 $\alpha'\in T'$ 和每个
$i\in\{1,\ldots,n\}$，映射列
$(\phi_{\alpha'}^\nu)^{-1}(z_{\alpha''i}^\nu)$ 收敛到 $z_{\alpha'i}'$。所以，对每个
$\alpha\in T$ 和每个 $i\in\{1,\ldots,n\}$，
\[
  d\bigl(\psi_\alpha^{-1}(z_{\alpha'i}'),z_{\alpha i}\bigr)
  =\lim_{\nu\to\infty}
  d\bigl((\chi_\alpha^\nu)^{-1}(z_{\alpha''i}^\nu),z_{\alpha i}\bigr).
\]
最后，若 $\alpha,\beta\in T$ 且 $\alpha E\beta$，则由
定义~\ref{def:5-5-1}，
\[
  E_{\alpha'}\bigl(\mathbf u';
    \psi_\alpha(B_\varepsilon(z_{\alpha\beta}))\bigr)
  =\lim_{\nu\to\infty}
  E_{\alpha''}\bigl(\mathbf u^\nu;
    \chi_\alpha^\nu(B_\varepsilon(z_{\alpha\beta}))\bigr).
\]
合并这些等式即得式~\eqref{eq:5-5-7}。最后对满足式~\eqref{eq:5-5-5} 的
$g$ 与 $\{\psi_\alpha\}$ 取下确界，便得到所需的上极限不等式。
\end{proof}
